\section{Keeping the first-row squares sharpens the alternating gate}
\label{sec:ordinary-alternating-87}

\paragraph{Status and scope.}
New root derivation, whole-read by the author and independently
reviewed in full by /root/imprimitive. Both mathematical reviews pass;
independent exact rational and Young-diagram checks also pass.
It uses the already proved exact content identities and full normalizer
lemma of HD-ALT-03--05. No new asymptotic character estimate is assumed.
The conclusion retains precisely the transitive ordinary-core and actual
inducing-field hypotheses of HD-ALT-04b.

\begin{lemma}[HD-ALT-06a: a sharper tail estimate]
\label{HD-ALT-06a}
For $n\ge15$ and $\lambda=(n-k,\mu)$ with $\mu\vdash k$ and
$1\le k<n/4$, put
\[
 T(n,k)=\binom n2-k(n-k+1),\qquad
 S_0(n,k)=\frac{(n-k-1)(n-k)(2n-2k-1)}6,
\]
\[
 V(n,k)=\frac{4\{T(n,k)^2+n(n-1)-3S_0(n,k)\}}
                    {n(n-1)(n-2)(n-3)}.
                                                               \tag{ALT87.1}
\]
For the ordinary irreducible alternating restriction of $S^\lambda$,
with degree $a$, its three-cycle and double-transposition ratios obey
\[
 \frac{1-u}{3}\ge\frac{k(n-1-k)}{(n-1)(n-2)},\qquad
                         v\le V(n,k).                         \tag{ALT87.2}
\]
Consequently every nonscalar finite-order linear normalizer has
projective support at least $a\delta(n,k)$, where
\[
 \delta(n,k)=\min\left\{\frac14,
       \frac{k(n-1-k)}{(n-1)(n-2)},\frac{1-V(n,k)}4\right\}.
                                                               \tag{ALT87.3}
\]
On the real interval $1\le k<n/4$, the second entry of this minimum
is increasing and $V(n,k)$ is strictly decreasing.
\end{lemma}

\begin{proof}
Write $t$ and $S$ for the sum of contents and their squares.
HD-ALT-05a gives $0<t\le T(n,k)$ and the first inequality in
(ALT87.2). The first row alone contributes exactly $S_0(n,k)$
to $S$, so $S\ge S_0(n,k)$. Keeping that contribution in the
exact identity (ALTQ.3) gives
\[
 v=\frac{4\{t^2+n(n-1)-3S\}}{n(n-1)(n-2)(n-3)}\le V(n,k).
\]
The long row prevents a self-conjugate label, so no split-constituent
degree adjustment is needed. The full normalizer conclusion follows
from HD-ALT-03a, with the harmless additional minimum $1/4$.

For monotonicity write $m=n-k$ and $N=T^2+n(n-1)-3S_0$.
The first derivative in (ALT87.2) is positive because
$n-1-2k>0$. Also
\[
 \frac{\partial N}{\partial k}
       =-2T(n+1-2k)+3m^2-3m+\frac12.                      \tag{ALT87.4}
\]
For $k\ge1$, the exact expression for $T$ gives
$T\ge n(n-2k-1)/2\ge n(n-2)/4$ on this interval.
Using $n+1-2k\ge n/2$ and $m\le n$ in (ALT87.4) gives
\[
 \frac{\partial N}{\partial k}
 \le-\frac{n^2(n-14)}4+\frac12<0\qquad(n\ge15).
\]
The denominator of $V$ is positive and independent of $k$.
This proves both monotonicity assertions.
\end{proof}

\begin{theorem}[HD-ALT-06b: the remaining row or column exceeds $87n/100$]
\label{HD-ALT-06b}
In the transitive ordinary alternating-core setting of HD-ALT-04b,
every factor label satisfies
\[
                       \max(\lambda_1,\lambda'_1)>87n/100.
                                                               \tag{ALT87.5}
\]
This allows every compatible outer operator, scalar phase, tensor
permutation and larger inducing character field in that earlier theorem.
More precisely, if a remaining label is written $(n-k,\mu)$ with
$k<n/4$ and total inducing degree $e$, then the explicit necessary
inequality is
\[
                          \delta(n,k)\le\frac1{10}+\frac3{5e}.
                                                               \tag{ALT87.6}
\]
The inequality does not assert the existence of a source whenever it holds.
\end{theorem}

\begin{proof}
By HD-ALT-04b, transpose if necessary so that $\lambda=(n-k,\mu)$
with $0<k<n/4$. Transitivity gives the same label up to transpose
on every factor. The ordinary factor is real, so the entire core has
a unique nondegenerate symmetric form line. The inducing image acts
on it by a nontrivial multiplier, exactly as in HD-ALT-04b: a trivial
multiplier would make its character real, contrary to the actual
inducing field containing the original imaginary quadratic field.

The actual intermediate-cover argument forces genus zero, at least
twelve genuine inertias and total moving rank $2e+6$. At least two
inertias have nontrivial multiplier and moving rank at least $e/2$.
The bound in HD-ALT-06a passes to tensor products by convolution;
nontrivial tensor cycles have the quarter bound, already included
in (ALT87.3). A genuine nonidentity scalar moves all dimensions.
Therefore
\[
                   2e+6\ge e+10e\delta(n,k),                \tag{ALT87.7}
\]
which is (ALT87.6). This remains a necessary constraint on the actual
period map, not a sufficiency assertion based on parameter counts.

Suppose (ALT87.5) fails, so $k\ge cn$ for $c=13/100$.
Put
\[
 A=1-2c+2c^2=\frac{3869}{5000},\qquad
 \delta_* =\frac{1-A^2}4=\frac{10030839}{100000000}>\frac1{10}.
                                                               \tag{ALT87.8}
\]
Direct polynomial simplification gives
\[
 \frac{cn(n-1-cn)}{(n-1)(n-2)}-c(1-c)
       =\frac{13(161n-174)}{10000(n-1)(n-2)}>0,             \tag{ALT87.9}
\]
and
\[
 A^2-V(n,cn)=
 \frac{13(1906518n^2-6812633n+4129618)}
              {25000000(n-1)(n-2)(n-3)}>0.                \tag{ALT87.10}
\]
For completeness, the numerator polynomial inside the parentheses,
at $n=15+x$, is
\[
              1906518x^2+50382907x+330906673>0.
\]
Since $c(1-c)=1131/10000>\delta_*$, monotonicity in the lemma
shows $\delta(n,k)>\delta_*$ whenever $k\ge cn$.
For $n\ge30$, this and (ALT87.7) force
\[
                  e<\frac{6}{10\delta_*-1}
                     =\frac{60000000}{30839}<2000.         \tag{ALT87.11}
\]
But $k\ge cn$ excludes the three low-degree shapes with tails one
or two. The precise published degree consequence of Rasala used
in HD-ALT-04--05 gives
$a\ge2n^3/27\ge2000$; see
\cite[proof of Theorem~9.7, p.~312]{HD:GKT2019}.
There is no restriction factor of two, since the long-row label is
not self-conjugate. Thus $e\ge a\ge2000$, a contradiction.

For $15\le n\le29$, monotonicity allows $k$ to be replaced by
$k_0=\lceil13n/100\rceil$. The complete exact endpoint table is
\[
\begin{array}{c|c|c@{\qquad}c|c|c}
n&k_0&\delta(n,k_0)&n&k_0&\delta(n,k_0)\\\hline
15&2&4001/32760&23&3&3989/35420\\
16&3&1091/7280&24&4&1875/14168\\
17&3&341/2380&25&4&19483/151800\\
18&3&1679/12240&26&4&573/4600\\
19&3&2039/15504&27&4&2831/23400\\
20&3&2447/19380&28&4&28891/245700\\
21&3&1453/11970&29&4&10865/95004\\
22&3&3419/29260&&&
\end{array}                                                     \tag{ALT87.12}
\]
Every displayed fraction is greater than $1/9$. Equation (ALT87.7)
therefore gives $e<54$. For $n=15,k=2$, the two possible Specht
degrees are $15(15-3)/2=90$ and $(15-1)(15-2)/2=91$.
For every other case the same published degree bound gives
$a\ge2n^3/27\ge250$; for $n=15$ and a larger $k$, this is also
the applicable bound. Hence $e\ge90$, the final contradiction.
\end{proof}

\paragraph{Strictness, exact checks and placement.}
The label $(19,3)$ at $n=22$ has row length greater than $6n/7$
but at most $87n/100$, so the reduction strictly narrows HD-ALT-05b.
The standalone exact certificate checks both polynomial identities,
the derivative identity, all fifteen displayed endpoint fractions,
and the degree ceilings. It does not enumerate all alternating groups
or decide the remaining still-longer shapes.

This is an additional restriction on actual finite-monodromy
weight-one variations and the condition-A supplier gate, cote16
scan050, with the precise EGA/SGA and representation-ring placement
in the common framework section. It supplies no cycle, no new
condition-I positivity argument, no Cayley--Hamilton inverse, and
no motivic $L$-value estimate. The existing exact detector constants
and the conditional status of Markman, Severn and032 are unaffected.
